% !TeX root = ../main.tex
\begin{abstract}
Accurate force and moment models are needed for flapping-wing simulation, control design, and flight-log interpretation, but physics-based models often leave systematic errors when applied to outdoor flight of a complete vehicle. This letter studies real-flight effective-wrench modeling for a bird-scale flapping-wing vehicle. We derive six-component body-frame effective wrench labels from onboard flight logs and align them with a DeLaurier-style simulator prior. Instead of replacing the prior with a black-box model, we formulate a theory-guided residual problem in which a causal neural model predicts the remaining discrepancy between the calibrated simulator prior and the log-derived effective wrench. On held-out whole-log tests, the calibrated prior has an overall RMSE of $5.97$, while a residual Transformer reduces the combined prediction error to $0.85$. In date-stratified five-fold whole-log evaluation, the reduction is stable, from $5.883\pm0.040$ to $0.827\pm0.027$ RMSE. Residual analysis shows that the remaining DeLaurier mismatch is not random: the main force-channel residuals contain repeatable wingbeat-phase, frequency, and flight-condition structure, and the neural residual removes most of this structured discrepancy. These results support theory-guided neural residual modeling as a practical route for bringing real-flight effects into flapping-wing simulation.
\end{abstract}
